\documentclass[12pt,a4paper]{article}
\usepackage[utf8]{inputenc}
\usepackage[spanish]{babel}
\usepackage{amsmath,amssymb,amsfonts}
\usepackage{geometry}
\usepackage{hyperref}
\geometry{margin=2.5cm}

\title{\textbf{Constantes Matemáticas en Cálculo Integral}}
\author{Compendio de Constantes Fundamentales}
\date{\today}

\begin{document}

\maketitle

\section*{Introducción}
El cálculo integral es la rama de las matemáticas orientada al cálculo de áreas, volúmenes, acumulaciones y la inversión de la derivación. Las constantes en cálculo integral surgen frecuentemente como valores de integrales definidas impropias o representaciones de la función Zeta de Riemann y constantes de integración universales.

\section{Logaritmo Natural de 2 ($\ln 2$)}

\subsection*{1. Significado, Símbolo y Explicación}
Denotado por $\ln 2$ o $\log_e 2$, se define de forma natural en cálculo integral como la integral definida de la función recíproca entre $1$ y $2$:
\begin{equation*}
\ln 2 = \int_{1}^{2} \frac{1}{x} \, dx
\end{equation*}
Aparece constantemente al evaluar integrales por partes, integrales por sustitución trigonométrica, al calcular el tiempo de vida media ($T_{1/2} = \frac{\ln 2}{\lambda}$) y en la suma de la serie armónica alternada ($\sum_{k=1}^\infty \frac{(-1)^{k+1}}{k} = \ln 2$).

\subsection*{2. Valor Típico en Ingeniería y Ciencia}
\begin{equation*}
\ln 2 \approx 0.6931 \quad \text{o} \quad 0.69314718
\end{equation*}

\subsection*{3. Valor Exacto a 15 Decimales}
\begin{equation*}
\ln 2 \approx 0.693147180559945
\end{equation*}

\subsection*{4. Valor Exacto a 100 Decimales}
\begin{align*}
\ln 2 \approx \; & 0.69314718055994530941723212145817656807550013436025525412068000949339... \\
& ...3621969694715605863326996418687542001
\end{align*}
\textbf{Margen de error:} Constante matemática pura (sin margen de error).

\vspace{0.8cm}
\hrule
\vspace{0.8cm}

\section{Constante de Apéry ($\zeta(3)$)}

\subsection*{1. Significado, Símbolo y Explicación}
Representada como $\zeta(3)$, es el valor de la función Zeta de Riemann evaluada en el entero $3$. Posee una representación directa como integral múltiple impropia:
\begin{equation*}
\zeta(3) = \sum_{n=1}^{\infty} \frac{1}{n^3} = \frac{1}{2} \int_{0}^{1} \int_{0}^{1} \int_{0}^{1} \frac{1}{1 - xyz} \, dx \, dy \, dz
\end{equation*}
En 1978, Roger Apéry demostró que es un número irracional. Surge de manera crucial en la integración de diagramas de Feynman en electrodinámica cuántica y física estadística.

\subsection*{2. Valor Típico en Ingeniería y Ciencia}
\begin{equation*}
\zeta(3) \approx 1.2021 \quad \text{o} \quad 1.20205690
\end{equation*}

\subsection*{3. Valor Exacto a 15 Decimales}
\begin{equation*}
\zeta(3) \approx 1.202056903159594
\end{equation*}

\subsection*{4. Valor Exacto a 100 Decimales}
\begin{align*}
\zeta(3) \approx \; & 1.20205690315959428539973816151144999076498629234049888179227155534183... \\
& ...82057863130901864558736093352581462
\end{align*}
\textbf{Margen de error:} Constante matemática pura (sin incertidumbre).

\end{document}
